\section{Discussion}
\label{sec:discussion}

We interpret the main findings and identify the calibration work needed for the next benchmark version.

\subsection{What the results say about current LLMs}

\paragraph{Classical methods still win on clean formal tasks.} PC + greedy achieved the highest directed F1 and the lowest SHD on the ladder. Under the v0 assumptions, this is consistent with theory: PC is asymptotically correct and the greedy orientation heuristic is strong enough for small graphs. Current LLM agents are not drop-in replacements for classical structure-learning algorithms in settings where the SCM lies in their identifiability regime.

\paragraph{The precision--recall asymmetry is the main behavioral signal.} PC under-commits (precision $66.9$, recall $32.8$); LLMs over-commit (precision $35$--$46$, recall $20$--$30$). Neither failure mode is uniformly preferable: under-committers leave useful edges unresolved, while over-committers add wrong orientations. A benchmark that reports only F1 collapses this distinction. ACDB's scoring contract keeps precision and recall visible.

\paragraph{Interventions sometimes hurt.} At L2 and L5, Sonnet 4.6's observational directed F1 exceeded its active directed F1 on identical seeds ($39.3\%$ vs $36.9\%$; $40.1\%$ vs $31.9\%$). Interventions are more informative than observations, so this is a policy diagnostic: current prompts do not reliably convert intervention samples into better orientation decisions. Edge-level revision analysis is a v1 target.

\paragraph{Statistical tools did not improve active discovery.} Giving LLMs correlation, partial-correlation, and independence-test tools did not raise directed F1 above raw data. The stats variants ran longer trajectories, intervened less, and submitted fewer directed edges; their $100\%$ efficiency reflects abstention, not skill. Tool access alone is insufficient if the model does not know when to stop testing and commit.

\subsection{What the results say about the benchmark}

\paragraph{Density leakage is an open calibration issue.} The random baseline's $23.6\%$ directed F1 on the current ladder is high enough that several GPT-5.4 settings sit near it. The density probe reduces the random floor to $16.9\%$ overall, turning the ladder issue into a concrete v1 calibration target rather than a vague caveat. The v0 results remain useful because all agents share paired seeds on the same ladder.

\paragraph{The scoring contract separates distinct failure modes.} The decomposition into skeleton, compelled-edge, DAG, SHD, and efficiency scores separated the CPDAG-ceiling question from the orientation-under-intervention question from the policy-quality question. The correlation-only probe would not be legible in an F1-only benchmark: its zero directed F1 looks like a failure, but the prompt-restricted reading --- the model withheld orientation when uncertain --- is only available because skeleton F1 is reported alongside.

\subsection{v1 and future variants}
\label{sec:future-scm-ladder}

The next benchmark version should keep the same SCM foundation but fix the density ladder, run the full LLM-correlation summary condition, increase seeds, and inspect intervention-driven revision behavior. Beyond that, the same generator/API can support causal inference with a known DAG, counterfactual queries, decision-making under causal uncertainty, transfer across SCM families, and RL for experimental design. Those are extensions of the scaffold, not claims made by this v0 run.

\subsection{Limitations}

Four limitations are central to interpreting the results. (1) v0 assumes causal sufficiency and linear--Gaussian SCMs. (2) The ladder has density-leakage artifacts at the dense end. (3) Eight seeds per level is adequate for coarse trends but under-powered for fine-grained ranking. (4) We tested two LLM backbones and one prompt design each, so LLM capability claims are conditional on those choices.
